\scanpage{89}

\[
\deleted{\unsure{f(\rho|f(\rho)}}\sum_1^{\unsure{k|n}}f(|\lambda_i|)\leq\sum_1^k f(\rho_i)\qquad\deleted{(k\leq\dim H)}\inserted{(k=1,\ldots)}\tag*{(18)}
\]

En particulier, \deleted{\unsure{pour que}} si \deleted{\unsure{$\dim$}} $\unsure{E|u}$ est de dimension infinie, si \inserted{la série} $\sum_{i=0}^{\infty}f(\rho_i)$ \deleted{\unsure{conv., alors|converge, alors}}\inserted{est convergente}, il en est de $=$ de $\sum_{\unsure{i|j}=0}^{\infty}f(|\lambda_i|)$, et on a

\[
\sum_{i=1}^{\infty}f(|\lambda_i|)\leq\sum_{i=0}^{\infty}f(\rho_i)\tag*{(19)}
\]

En particulier, faisant $f(\unsure{r|a})=r^p$ $(p>0)$, on trouve

\uline{Corollaire 2} Soit $u$ un op. compact dans un espace de Hilbert $E$, et soit $p>0$. Alors, pour tout entier $k>0$

\[
\boxed{\sum_{i=1}^k|\lambda_i|^p\leq\sum_{\unsure{j|i}=1}^k\rho_i^p\deleted{=\|u\|_p}}\qquad\deleted{(k\leq\dim H)}\tag*{(20)}
\]

\deleted{en particulier, si $E$ est de dimension finie $n$}

\[
\boxed{\sum_i|\lambda_i|^p\leq\sum_i\rho_i^p=\|u\|_p}\tag*{(21)}
\]

(où la somme est prise sur tous les indices $i$) En particulier, faisant $p=1$ et $p=2$, et appliquant le N° 1, th. 2 et th. 3, formules (15) et (20), on trouve

\uline{Corollaire 3} Soit $u$ un op. \deleted{\unsure{dont}} de Fredholm (resp.t un op. de \deleted{\unsure{Hilb.|Hilb}} H.-S.) dans \deleted{\unsure{l'aire}}\inserted{un espace} de Hilbert. \deleted{\unsure{Et}} Alors \deleted{\unsure{une}} la suite des valeurs propres est sommable (resp.t de carré sommable) et on a

\[
\boxed{\sum|\lambda_i|\leq\|u\|_1}\tag*{(21)}
\]

\[
\boxed{\left(\sum|\lambda_i|^2\right)^{1/2}\leq\|u\|_2}\tag*{(22)}
\]

Signalons aussi le

% slanted
\margin{petits caractères}

\uline{Corollaire 4} Soit $u$ un op. \deleted{\unsure{dans}}\inserted{lin.} dans un espace de Hilbert de dimension finie $n$, \deleted{Alors,} soit $f(s_1,\ldots s_n)$ \deleted{\unsure{est}} une fonction de $n$ arguments $s_i>0$ telle que $f(e^{t_1},\ldots e^{t_n})$ soit convexe (pas forcément croissante). Si $f$ est en fait définie et continue pour \deleted{\unsure{les}} $s_i\geq0$, ou si $u$ est inversible, on a

\[
\boxed{f(|\lambda_1|,\ldots|\lambda_n|)\leq f(\rho_1,\ldots,\rho_n)}.\tag*{\textcircled{23}}
\]

\deleted{\unsure{(23)|Alors}} En effet, on a, en plus des inégalités $|\lambda_1\ldots\lambda_{\unsure{n|k}}|\leq\rho_1\ldots\rho_n$ $(k\leq n)$, l'\uline{égalité}

% archive sheet number: 90
